\documentclass[11pt,a4paper]{article}
\usepackage[utf8]{inputenc}
\usepackage[T1]{fontenc}
\usepackage[spanish,es-minimal,es-noshorthands]{babel}
\usepackage{lmodern}
\usepackage[sfdefault,lining]{FiraSans}
\usepackage[a4paper,top=2.3cm,bottom=2.2cm,left=1.8cm,right=1.8cm,headheight=1.3cm,headsep=0.4cm,footskip=1cm]{geometry}
\usepackage{amsmath,amssymb}
\usepackage{icomma}
\usepackage{cancel}
\usepackage{graphicx}
\usepackage[table]{xcolor}
\usepackage{tikz}
\usepackage{fontawesome5}
\usepackage{ccicons}
\usepackage{fancyhdr}
\usepackage{lastpage}
\usepackage{tasks}
\usepackage{needspace}
\usepackage[most]{tcolorbox}
\usepackage{array}
\usepackage[hidelinks]{hyperref}
\AtBeginDocument{\hypersetup{pdftitle={Soluciones de la ficha de la unidad 1: Números naturales. 1º ESO},pdfsubject={Matemáticas 1º ESO},pdfkeywords={números naturales, aproximación, redondeo, truncamiento, operaciones combinadas, potencias, potencias de 10, raíces cuadradas}}}

\definecolor{navy}{HTML}{1B2748}
\definecolor{accent}{HTML}{E8603A}
\definecolor{teal}{HTML}{1F9E90}
\definecolor{soft}{HTML}{EEF1F8}
\definecolor{muted}{HTML}{5D6685}

\setlength{\parindent}{0pt}
\setlength{\parskip}{3pt}
\settasks{label=\alph*), label-format=\bfseries\color{navy}, label-width=1.3em, item-indent=1.8em, column-sep=0.6em, after-item-skip=2pt, before-skip=3pt, after-skip=4pt}

% ---------- capçalera i peu ----------
\newcommand{\llicencia}{\href{https://creativecommons.org/licenses/by-nc-sa/4.0/deed.es}{\ccbyncsa}\ \ Licencia Creative Commons Atribución-NoComercial-CompartirIgual 4.0 Internacional (CC BY-NC-SA 4.0)}
\pagestyle{fancy}
\fancyhf{}
\fancyhead[L]{\raisebox{-0.35\height}{\includegraphics[height=0.95cm]{logo-gva.png}}}
\fancyhead[C]{\small\color{muted}Soluciones de la ficha · Unidad 1: Números naturales · 1\textsuperscript{o} ESO}
\fancyhead[R]{\raisebox{-0.35\height}{\includegraphics[height=0.95cm]{logo-vermellar.png}}}
\fancyfoot[L]{\scriptsize\color{muted}\llicencia}
\fancyfoot[R]{\small\color{muted}\thepage}
\renewcommand{\headrulewidth}{0.4pt}
\renewcommand{\headrule}{\color{muted!40}\hrule height \headrulewidth}
\fancypagestyle{primera}{\fancyhf{}\renewcommand{\headrulewidth}{0pt}
  \fancyfoot[L]{\scriptsize\color{muted}\llicencia}
  \fancyfoot[R]{\small\color{muted}\thepage}}

% ---------- nivells i exercicis ----------
\newcommand{\estrelles}[1]{{\color{accent}\ifcase#1\or\faStar\faStar[regular]\faStar[regular]\or\faStar\faStar\faStar[regular]\or\faStar\faStar\faStar\fi}}
\newcounter{exer}
\newcommand{\ex}[1]{\par\needspace{4\baselineskip}\medskip\stepcounter{exer}\noindent{\color{navy}\bfseries\theexer.}\ \ {\small\estrelles{#1}}\ \ }

\newcommand{\seccio}[2]{\par\needspace{10\baselineskip}\bigskip
  \begin{tcolorbox}[colback=navy,colframe=navy,coltext=white,arc=2mm,boxrule=0pt,left=3mm,right=3mm,top=1.2mm,bottom=1.2mm]
  {\large\bfseries #1\quad #2}
  \end{tcolorbox}\vspace{-2pt}}

\newcolumntype{C}[1]{>{\centering\arraybackslash}p{#1}}

% quadre de recordatori (només a la fitxa d'exercicis)
\newcommand{\recorda}[1]{\par\smallskip\begin{tcolorbox}[colback=accent!7,colframe=accent!55,arc=2mm,boxrule=0.5pt,left=2.5mm,right=2.5mm,top=1mm,bottom=1mm]\small{\color{accent}\faLightbulb[regular]}\ \textbf{Recuerda:}\ #1\end{tcolorbox}}
\newcommand{\res}[1]{\boxed{#1}}

% =====================================================================
% Capçalera FentMates de les fitxes: un color per a cada curs.
% Els colors són els de la pàgina inicial (LEVELS a assets/js/materials.js);
% el logo ve de _eines/logo/pdf/ (python3 _eines/logo/genera.py).
%
% Ús, al final del preàmbul de la unitat:
%   % =====================================================================
% Capçalera FentMates de les fitxes: un color per a cada curs.
% Els colors són els de la pàgina inicial (LEVELS a assets/js/materials.js);
% el logo ve de _eines/logo/pdf/ (python3 _eines/logo/genera.py).
%
% Ús, al final del preàmbul de la unitat:
%   % =====================================================================
% Capçalera FentMates de les fitxes: un color per a cada curs.
% Els colors són els de la pàgina inicial (LEVELS a assets/js/materials.js);
% el logo ve de _eines/logo/pdf/ (python3 _eines/logo/genera.py).
%
% Ús, al final del preàmbul de la unitat:
%   \input{../comu/fentmates.tex}
%   \fmcurs{1eso}
% i, al començament del document:
%   \capcalera{Matemàtiques}{1r ESO}{Fitxa d'exercicis · Unitat 2}{Divisibilitat}
% Opcional abans de \input: \def\fmneed{14} (espai mínim abans de cada secció).
% =====================================================================
\makeatletter
\def\fm@logodir{../../logo/pdf}
\providecommand{\fmneed}{10}
\def\fm@cursnom{}

% clau del curs, color i color del text damunt del color
\newcommand{\fm@curs}[3]{\@namedef{fm@color@#1}{#2}\@namedef{fm@text@#1}{#3}}
\fm@curs{1eso}{E6B84A}{navy}
\fm@curs{2eso}{3B82F6}{white}
\fm@curs{3eso}{FF7A50}{white}
\fm@curs{4eso-a}{A78BFA}{white}
\fm@curs{4eso-b}{6D28D9}{white}
\fm@curs{1batx-mi}{9B1C1C}{white}
\fm@curs{1batx-ccss}{EF6B6B}{white}
\fm@curs{2batx-mii}{1E6B47}{white}
\fm@curs{2batx-ccss}{2BB3A3}{white}

\newcommand{\fmcurs}[1]{%
  \@ifundefined{fm@color@#1}{\PackageError{fentmates}{Curs desconegut: #1}{}}{}%
  \edef\fm@tmp{\noexpand\definecolor{curs}{HTML}{\@nameuse{fm@color@#1}}%
               \noexpand\colorlet{curstext}{\@nameuse{fm@text@#1}}}\fm@tmp
  \colorlet{cursfosc}{curs!72!black}%
  \def\fm@logo{\fm@logodir/vertical-#1.pdf}%
  \def\fm@marca{\fm@logodir/marca-#1.pdf}}

% ---------- capçalera de la primera pàgina ----------
% #1 matèria (petit), #2 curs, #3 tipus de document i unitat, #4 títol
\newcommand{\capcalera}[4]{%
  \gdef\fm@cursnom{#2}%
  \par\noindent
  \begin{tcolorbox}[colback=curs,colframe=curs,boxrule=0pt,arc=2.2mm,width=3.5cm,height=2.15cm,
      valign=center,left=3mm,right=2mm,top=1mm,bottom=1mm,nobeforeafter,box align=center]
    \color{curstext}{\footnotesize #1}\par\vspace{1pt}{\fontsize{19}{21}\selectfont\bfseries #2\par}%
  \end{tcolorbox}\hspace{5mm}%
  \begin{minipage}[c]{\dimexpr\linewidth-3.5cm-5mm-2.9cm\relax}
    {\color{cursfosc}\small\bfseries\MakeUppercase{#3}}\par\vspace{2pt}
    {\color{navy}\fontsize{20}{23}\selectfont\bfseries #4\par}
  \end{minipage}\hfill
  \begin{minipage}[c]{2.6cm}\raggedleft\includegraphics[width=2.6cm]{\fm@logo}\end{minipage}\par
  \vspace{2pt}}

% ---------- barres de secció del color del curs ----------
\renewcommand{\seccio}[2]{\par\needspace{\fmneed\baselineskip}\bigskip
  \begin{tcolorbox}[colback=curs,colframe=curs,coltext=curstext,arc=2mm,boxrule=0pt,left=3mm,right=3mm,top=1.2mm,bottom=1.2mm]
  {\large\bfseries #1\quad #2}
  \end{tcolorbox}\vspace{-2pt}}

% ---------- estrelles de dificultat del color del curs (a Batxillerat no se'n posen) ----------
\@ifundefined{estrelles}{}{%
  \renewcommand{\estrelles}[1]{{\color{curs}\ifcase#1\or\faStar\faStar[regular]\faStar[regular]\or\faStar\faStar\faStar[regular]\or\faStar\faStar\faStar\fi}}}

% ---------- capçalera i peu de les pàgines ----------
% primera pàgina: logos institucionals al peu
\fancypagestyle{primera}{\fancyhf{}\renewcommand{\headrulewidth}{0pt}
  \fancyfoot[L]{\raisebox{-0.32\height}{\includegraphics[height=0.78cm]{logo-gva.png}}\hspace{5mm}%
    \raisebox{-0.32\height}{\includegraphics[height=0.78cm]{logo-vermellar.png}}}
  \fancyfoot[R]{\scriptsize\color{muted}\href{https://creativecommons.org/licenses/by-nc-sa/4.0/deed.ca}{\ccbyncsa}\ \ CC BY-NC-SA 4.0 · fentmates.org\qquad\small\thepage}}
% la resta: el símbol a l'esquerra i el curs a la dreta, amb el títol al centre com abans
\fancyhead[L]{\raisebox{-0.3\height}{\includegraphics[height=0.62cm]{\fm@marca}}}
\fancyhead[R]{\small\bfseries\color{cursfosc}\fm@cursnom}
\renewcommand{\headrule}{\color{curs!45}\hrule height \headrulewidth}
\makeatother

%   \fmcurs{1eso}
% i, al començament del document:
%   \capcalera{Matemàtiques}{1r ESO}{Fitxa d'exercicis · Unitat 2}{Divisibilitat}
% Opcional abans de \input: \def\fmneed{14} (espai mínim abans de cada secció).
% =====================================================================
\makeatletter
\def\fm@logodir{../../logo/pdf}
\providecommand{\fmneed}{10}
\def\fm@cursnom{}

% clau del curs, color i color del text damunt del color
\newcommand{\fm@curs}[3]{\@namedef{fm@color@#1}{#2}\@namedef{fm@text@#1}{#3}}
\fm@curs{1eso}{E6B84A}{navy}
\fm@curs{2eso}{3B82F6}{white}
\fm@curs{3eso}{FF7A50}{white}
\fm@curs{4eso-a}{A78BFA}{white}
\fm@curs{4eso-b}{6D28D9}{white}
\fm@curs{1batx-mi}{9B1C1C}{white}
\fm@curs{1batx-ccss}{EF6B6B}{white}
\fm@curs{2batx-mii}{1E6B47}{white}
\fm@curs{2batx-ccss}{2BB3A3}{white}

\newcommand{\fmcurs}[1]{%
  \@ifundefined{fm@color@#1}{\PackageError{fentmates}{Curs desconegut: #1}{}}{}%
  \edef\fm@tmp{\noexpand\definecolor{curs}{HTML}{\@nameuse{fm@color@#1}}%
               \noexpand\colorlet{curstext}{\@nameuse{fm@text@#1}}}\fm@tmp
  \colorlet{cursfosc}{curs!72!black}%
  \def\fm@logo{\fm@logodir/vertical-#1.pdf}%
  \def\fm@marca{\fm@logodir/marca-#1.pdf}}

% ---------- capçalera de la primera pàgina ----------
% #1 matèria (petit), #2 curs, #3 tipus de document i unitat, #4 títol
\newcommand{\capcalera}[4]{%
  \gdef\fm@cursnom{#2}%
  \par\noindent
  \begin{tcolorbox}[colback=curs,colframe=curs,boxrule=0pt,arc=2.2mm,width=3.5cm,height=2.15cm,
      valign=center,left=3mm,right=2mm,top=1mm,bottom=1mm,nobeforeafter,box align=center]
    \color{curstext}{\footnotesize #1}\par\vspace{1pt}{\fontsize{19}{21}\selectfont\bfseries #2\par}%
  \end{tcolorbox}\hspace{5mm}%
  \begin{minipage}[c]{\dimexpr\linewidth-3.5cm-5mm-2.9cm\relax}
    {\color{cursfosc}\small\bfseries\MakeUppercase{#3}}\par\vspace{2pt}
    {\color{navy}\fontsize{20}{23}\selectfont\bfseries #4\par}
  \end{minipage}\hfill
  \begin{minipage}[c]{2.6cm}\raggedleft\includegraphics[width=2.6cm]{\fm@logo}\end{minipage}\par
  \vspace{2pt}}

% ---------- barres de secció del color del curs ----------
\renewcommand{\seccio}[2]{\par\needspace{\fmneed\baselineskip}\bigskip
  \begin{tcolorbox}[colback=curs,colframe=curs,coltext=curstext,arc=2mm,boxrule=0pt,left=3mm,right=3mm,top=1.2mm,bottom=1.2mm]
  {\large\bfseries #1\quad #2}
  \end{tcolorbox}\vspace{-2pt}}

% ---------- estrelles de dificultat del color del curs (a Batxillerat no se'n posen) ----------
\@ifundefined{estrelles}{}{%
  \renewcommand{\estrelles}[1]{{\color{curs}\ifcase#1\or\faStar\faStar[regular]\faStar[regular]\or\faStar\faStar\faStar[regular]\or\faStar\faStar\faStar\fi}}}

% ---------- capçalera i peu de les pàgines ----------
% primera pàgina: logos institucionals al peu
\fancypagestyle{primera}{\fancyhf{}\renewcommand{\headrulewidth}{0pt}
  \fancyfoot[L]{\raisebox{-0.32\height}{\includegraphics[height=0.78cm]{logo-gva.png}}\hspace{5mm}%
    \raisebox{-0.32\height}{\includegraphics[height=0.78cm]{logo-vermellar.png}}}
  \fancyfoot[R]{\scriptsize\color{muted}\href{https://creativecommons.org/licenses/by-nc-sa/4.0/deed.ca}{\ccbyncsa}\ \ CC BY-NC-SA 4.0 · fentmates.org\qquad\small\thepage}}
% la resta: el símbol a l'esquerra i el curs a la dreta, amb el títol al centre com abans
\fancyhead[L]{\raisebox{-0.3\height}{\includegraphics[height=0.62cm]{\fm@marca}}}
\fancyhead[R]{\small\bfseries\color{cursfosc}\fm@cursnom}
\renewcommand{\headrule}{\color{curs!45}\hrule height \headrulewidth}
\makeatother

%   \fmcurs{1eso}
% i, al començament del document:
%   \capcalera{Matemàtiques}{1r ESO}{Fitxa d'exercicis · Unitat 2}{Divisibilitat}
% Opcional abans de \input: \def\fmneed{14} (espai mínim abans de cada secció).
% =====================================================================
\makeatletter
\def\fm@logodir{../../logo/pdf}
\providecommand{\fmneed}{10}
\def\fm@cursnom{}

% clau del curs, color i color del text damunt del color
\newcommand{\fm@curs}[3]{\@namedef{fm@color@#1}{#2}\@namedef{fm@text@#1}{#3}}
\fm@curs{1eso}{E6B84A}{navy}
\fm@curs{2eso}{3B82F6}{white}
\fm@curs{3eso}{FF7A50}{white}
\fm@curs{4eso-a}{A78BFA}{white}
\fm@curs{4eso-b}{6D28D9}{white}
\fm@curs{1batx-mi}{9B1C1C}{white}
\fm@curs{1batx-ccss}{EF6B6B}{white}
\fm@curs{2batx-mii}{1E6B47}{white}
\fm@curs{2batx-ccss}{2BB3A3}{white}

\newcommand{\fmcurs}[1]{%
  \@ifundefined{fm@color@#1}{\PackageError{fentmates}{Curs desconegut: #1}{}}{}%
  \edef\fm@tmp{\noexpand\definecolor{curs}{HTML}{\@nameuse{fm@color@#1}}%
               \noexpand\colorlet{curstext}{\@nameuse{fm@text@#1}}}\fm@tmp
  \colorlet{cursfosc}{curs!72!black}%
  \def\fm@logo{\fm@logodir/vertical-#1.pdf}%
  \def\fm@marca{\fm@logodir/marca-#1.pdf}}

% ---------- capçalera de la primera pàgina ----------
% #1 matèria (petit), #2 curs, #3 tipus de document i unitat, #4 títol
\newcommand{\capcalera}[4]{%
  \gdef\fm@cursnom{#2}%
  \par\noindent
  \begin{tcolorbox}[colback=curs,colframe=curs,boxrule=0pt,arc=2.2mm,width=3.5cm,height=2.15cm,
      valign=center,left=3mm,right=2mm,top=1mm,bottom=1mm,nobeforeafter,box align=center]
    \color{curstext}{\footnotesize #1}\par\vspace{1pt}{\fontsize{19}{21}\selectfont\bfseries #2\par}%
  \end{tcolorbox}\hspace{5mm}%
  \begin{minipage}[c]{\dimexpr\linewidth-3.5cm-5mm-2.9cm\relax}
    {\color{cursfosc}\small\bfseries\MakeUppercase{#3}}\par\vspace{2pt}
    {\color{navy}\fontsize{20}{23}\selectfont\bfseries #4\par}
  \end{minipage}\hfill
  \begin{minipage}[c]{2.6cm}\raggedleft\includegraphics[width=2.6cm]{\fm@logo}\end{minipage}\par
  \vspace{2pt}}

% ---------- barres de secció del color del curs ----------
\renewcommand{\seccio}[2]{\par\needspace{\fmneed\baselineskip}\bigskip
  \begin{tcolorbox}[colback=curs,colframe=curs,coltext=curstext,arc=2mm,boxrule=0pt,left=3mm,right=3mm,top=1.2mm,bottom=1.2mm]
  {\large\bfseries #1\quad #2}
  \end{tcolorbox}\vspace{-2pt}}

% ---------- estrelles de dificultat del color del curs (a Batxillerat no se'n posen) ----------
\@ifundefined{estrelles}{}{%
  \renewcommand{\estrelles}[1]{{\color{curs}\ifcase#1\or\faStar\faStar[regular]\faStar[regular]\or\faStar\faStar\faStar[regular]\or\faStar\faStar\faStar\fi}}}

% ---------- capçalera i peu de les pàgines ----------
% primera pàgina: logos institucionals al peu
\fancypagestyle{primera}{\fancyhf{}\renewcommand{\headrulewidth}{0pt}
  \fancyfoot[L]{\raisebox{-0.32\height}{\includegraphics[height=0.78cm]{logo-gva.png}}\hspace{5mm}%
    \raisebox{-0.32\height}{\includegraphics[height=0.78cm]{logo-vermellar.png}}}
  \fancyfoot[R]{\scriptsize\color{muted}\href{https://creativecommons.org/licenses/by-nc-sa/4.0/deed.ca}{\ccbyncsa}\ \ CC BY-NC-SA 4.0 · fentmates.org\qquad\small\thepage}}
% la resta: el símbol a l'esquerra i el curs a la dreta, amb el títol al centre com abans
\fancyhead[L]{\raisebox{-0.3\height}{\includegraphics[height=0.62cm]{\fm@marca}}}
\fancyhead[R]{\small\bfseries\color{cursfosc}\fm@cursnom}
\renewcommand{\headrule}{\color{curs!45}\hrule height \headrulewidth}
\makeatother

\fmcurs{1eso}

\begin{document}
\thispagestyle{primera}

\capcalera{Matemáticas}{1\textsuperscript{o} ESO}{Soluciones de la ficha · Unidad 1}{Números naturales}
\vspace{8pt}

\begin{tcolorbox}[colback=soft,colframe=soft,arc=2mm,left=3mm,right=3mm,top=1.5mm,bottom=1.5mm]
\small Intenta resolver cada ejercicio antes de mirar su solución. Si no te sale el mismo resultado, repasa los pasos o consulta los apuntes interactivos de la unidad.
\end{tcolorbox}

% =====================================================
\seccio{1}{Los números naturales}

\ex{1} \begin{tasks}(4)
  \task $23\,407$
  \task $105\,060$
  \task $2\,300\,000$
  \task $7\,050\,009$
\end{tasks}

\ex{1} \begin{tasks}(2)
  \task Cuatro mil noventa.
  \task Quince mil seiscientos.
  \task Trescientos ocho mil doce.
  \task Dos millones cuarenta mil quinientos.
\end{tasks}

\ex{1} \begin{tasks}(6)
  \task $700$
  \task $70\,000$
  \task $7$
  \task $7\,000\,000$
  \task $70$
  \task $70\,000$
\end{tasks}

\ex{2} $4599<45\,099<45\,909<45\,990<49\,059$ \ (el de cuatro cifras es el más pequeño; los demás se comparan cifra a cifra de izquierda a derecha).

\ex{2} \begin{tasks}(2)
  \task $999$ \ y \ $1001$
  \task $39\,998$ \ y \ $40\,000$
  \task $99\,999$ \ y \ $100\,001$
  \task $2\,009\,998$ \ y \ $2\,010\,000$
\end{tasks}

\ex{2} \begin{tasks}(3)
  \task $\res{345}$ decenas
  \task $\res{720}$ centenas
  \task $\res{1000}$ millares
\end{tasks}

\ex{3} \begin{tasks}(1)
  \task $\res{96\,420}$ (cifras de mayor a menor).
  \task $\res{20\,469}$ (empezamos por la cifra más pequeña que no es 0).
  \task Tiene que acabar en 9, la única cifra impar: $\res{64\,209}$.
\end{tasks}

\ex{3} Del $500$ al $599$: $\res{100}$ números. \ Del $1000$ al $9999$: $9999-999=\res{9000}$ números.


% =====================================================
\seccio{2}{Aproximación: truncamiento y redondeo}

\ex{1} \begin{tasks}(3)
  \task $1250$ \ y \ $1200$
  \task $32\,040$ \ y \ $32\,000$
  \task $57\,620$ \ y \ $57\,600$
  \task $12\,350$ \ y \ $12\,300$
  \task $3460$ \ y \ $3400$
  \task $23\,465\,240$ \ y \ $23\,465\,200$
\end{tasks}

\ex{1} \begin{tasks}(3)
  \task $1260$ \ y \ $1300$
  \task $32\,040$ \ y \ $32\,000$
  \task $57\,620$ \ y \ $57\,600$
  \task $12\,350$ \ y \ $12\,400$
  \task $3460$ \ y \ $3500$
  \task $23\,465\,240$ \ y \ $23\,465\,200$
\end{tasks}\par No coinciden en a) (decenas y centenas), d) (centenas) y e) (centenas). El \textbf{redondeo} siempre da el número más próximo.

\ex{2} \par\smallskip{\renewcommand{\arraystretch}{1.5}\centering\begin{tabular}{|C{3.6cm}|C{3.6cm}|C{3.6cm}|}\hline\rowcolor{soft}\textbf{Número} & \textbf{Truncamiento} & \textbf{Redondeo}\\\hline $45\,678$ & $45\,000$ & $46\,000$\\\hline $129\,350$ & $129\,000$ & $129\,000$\\\hline $9499$ & $9000$ & $9000$\\\hline $70\,501$ & $70\,000$ & $71\,000$\\\hline $999\,600$ & $999\,000$ & $1\,000\,000$\\\hline\end{tabular}\par}

\ex{2} \begin{tasks}(4)
  \task $4000$
  \task $60\,000$
  \task $200\,000$
  \task $3\,000\,000$
\end{tasks}

\ex{2} Decenas: $\res{1350\text{ €}}$, error $1350-1349=1$ €. \ Centenas: $\res{1300\text{ €}}$, error $1349-1300=49$ €.

\ex{3} Del $\res{8250}$ (sube) hasta el $\res{8349}$ (se queda igual).

\ex{3} El primero ha \textbf{redondeado} (error $348$) y el segundo ha \textbf{truncado} (error $652$). Se equivoca menos el primero.


% =====================================================
\seccio{3}{Operaciones con números naturales}

\ex{1} \begin{tasks}(4)
  \task $16\,479$
  \task $12\,614$
  \task $18\,252$
  \task $920\,000$
\end{tasks}

\ex{1} \begin{tasks}(2)
  \task $q=583,\ r=0$ \ (exacta)
  \task $q=384,\ r=20$: \ $24\cdot384+20=9236$
  \task $q=416,\ r=0$ \ (exacta)
  \task $q=111,\ r=7$: \ $63\cdot111+7=7000$
\end{tasks}

\ex{2} \par\smallskip{\renewcommand{\arraystretch}{1.4}\centering\begin{tabular}{|C{2.6cm}|C{2.6cm}|C{2.6cm}|C{2.6cm}|}\hline\rowcolor{soft}\textbf{Dividendo} & \textbf{Divisor} & \textbf{Cociente} & \textbf{Resto}\\\hline $427$ & $12$ & $35$ & $7$\\\hline $500$ & $12$ & $41$ & $8$\\\hline $1000$ & $27$ & $37$ & $1$\\\hline $474$ & $25$ & $18$ & $24$\\\hline\end{tabular}\par}\smallskip Para el divisor: $(500-8):41=492:41=12$.

\ex{2} \begin{tasks}(3)
  \task $(25\cdot4)\cdot17=\res{1700}$
  \task $100+100=\res{200}$
  \task $13\cdot100-13=\res{1287}$
  \task $\res{1000}$
  \task $36\cdot(7+3)=\res{360}$
  \task $45\cdot100+45\cdot2=\res{4590}$
\end{tasks}

\ex{1} \begin{tasks}(4)
  \task $5+15-12=\res{8}$
  \task $28-6-10=\res{12}$
  \task $4+7=\res{11}$
  \task $27-20=\res{7}$
  \task $30-20+5=\res{15}$
  \task $8-8=\res{0}$
  \task $27-8=\res{19}$
  \task $50-36=\res{14}$
\end{tasks}

\ex{2} \begin{tasks}(1)
  \task $24+9-15=\res{18}$
  \task $45-30-6\cdot2=45-30-12=\res{3}$
  \task $2\cdot10+10-18=\res{12}$
  \task $5\cdot[2+15-7]+7=5\cdot10+7=\res{57}$
  \task $[3+8]\cdot2-8=22-8=\res{14}$
  \task $25-12-[20-18]=13-2=\res{11}$
\end{tasks}

\ex{3} \begin{tasks}(1)
  \task $[50-42]\cdot3-2\cdot[12-4]=24-16=\res{8}$
  \task $100-3\cdot[20-2\cdot3]:7=100-3\cdot14:7=100-42:7=100-6=\res{94}$
\end{tasks}

\ex{3} \begin{tasks}(3)
  \task $(6+4)\cdot5-3=47$
  \task $6+4\cdot(5-3)=14$
  \task $(6+4)\cdot(5-3)=20$
\end{tasks}


% =====================================================
\seccio{4}{Problemas con números naturales}

\ex{1} $612:36=\res{17\text{ €}}$ cada balón.

\ex{1} Trabaja $31-6=25$ días y hace $13\cdot12=156$ pasteles al día: \ $156\cdot25=\res{3900\text{ pasteles}}$.

\ex{1} De cada caja sobreviven $758-27=731$: \ $731\cdot45=\res{32\,895\text{ alevines}}$.

\ex{2} $(34+12)\cdot15=46\cdot15=690$ kg; \ $690\cdot2=\res{1380\text{ €}}$.

\ex{2} $43\cdot7\cdot147\cdot15=301\cdot147\cdot15=44\,247\cdot15=\res{663\,705\text{ €}}$.

\ex{2} Alumnos: $28+27+29+26=110$; personas: $118$. \ $118:55\to q=2,\ r=8$: hacen falta \res{3\text{ autobuses}}.\\ $3\cdot330=990$ €; \ $990:110=\res{9\text{ €}}$ por alumno.

\ex{2} Rotas: $23\,400:3=7800$. \ Enteras: $23\,400-7800=15\,600$. \ $15\,600:12=\res{1300\text{ cajas}}$.

\ex{2} Precio de un kilo: $9100:1300=7$ €. \ $1750\cdot7=\res{12\,250\text{ €}}$.

\ex{3} De cada 5 pago 4: de los 20 he pagado $20:5\cdot4=16$. \ $64:16=\res{4\text{ €}}$ cada rotulador.

\ex{3} Asientos: $24\cdot18=432$; vacíos: $432-350=\res{82}$. \ Recaudado: $350\cdot7=\res{2450\text{ €}}$.\\ Lleno: $432\cdot7=3024$ €; \ $3024-2450=\res{574\text{ €}}$ más.

\ex{3} Faltan $1380-300=1080$ €. \ $1080:90=\res{12\text{ cuotas}}$ (un año).


% =====================================================
\seccio{5}{Potencias}

\ex{1} \begin{tasks}(2)
  \task $3\cdot3=9$
  \task $5\cdot5\cdot5=125$
  \task $7$
  \task $2\cdot2\cdot2\cdot2\cdot2=32$
  \task $9\cdot9=81$
  \task $6\cdot6\cdot6=216$
  \task $0\cdot0\cdot0\cdot0\cdot0=0$
  \task $1\cdot1\cdot1\cdot1\cdot1\cdot1\cdot1=1$
  \task $4\cdot4\cdot4\cdot4=256$
  \task $11\cdot11=121$
\end{tasks}

\ex{1} \begin{tasks}(2)
  \task $7^{4}$: siete elevado a cuatro
  \task $12^{2}$: doce al cuadrado
  \task $2^{6}$: dos elevado a seis
  \task $8^{3}$: ocho al cubo
\end{tasks}

\ex{2} \begin{tasks}(3)
  \task $7^{2}$
  \task $5^{3}$
  \task $2^{5}$
  \task $10^{3}$
  \task $9^{2}=3^{4}$
  \task $8^{2}=4^{3}=2^{6}$
\end{tasks}

\ex{2} \begin{tasks}(4)
  \task $32>25$
  \task $81>64$
  \task $16=16$
  \task $100<1024$
\end{tasks}

\ex{2} \par\smallskip{\renewcommand{\arraystretch}{1.4}\centering\begin{tabular}{|C{1.6cm}|C{1.6cm}|C{1.6cm}|C{1.6cm}|C{1.6cm}|C{1.6cm}|}\hline\rowcolor{soft}$n$ & $4$ & $7$ & $11$ & $13$ & $20$\\\hline $n^{2}$ & $16$ & $49$ & $121$ & $169$ & $400$\\\hline $n^{3}$ & $64$ & $343$ & $1331$ & $2197$ & $8000$\\\hline\end{tabular}\par}

\ex{3} Flores: $3\cdot3\cdot3\cdot3=3^{4}=\res{81}$. \ Llaves: $6\cdot6\cdot6=6^{3}=\res{216}$.

\ex{3} 1 hora: $2$; \ 2 horas: $2^{2}=4$; \ 10 horas: $2^{10}=\res{1024\text{ bacterias}}$.


% =====================================================
\seccio{6}{Potencias de base 10}

\ex{1} \begin{tasks}(3)
  \task $1000$
  \task $1\,000\,000$
  \task $100\,000\,000$
  \task $10^{4}$
  \task $10^{5}$
  \task $10^{1}$
\end{tasks}

\ex{1} \begin{tasks}(4)
  \task $6\cdot10^{6}$
  \task $47\cdot10^{3}$
  \task $35\cdot10^{7}$
  \task $208\cdot10^{4}$
\end{tasks}

\ex{1} \begin{tasks}(4)
  \task $900\,000$
  \task $130\,000$
  \task $207\,000$
  \task $500\,000\,000$
\end{tasks}

\ex{2} \begin{tasks}(2)
  \task $6\cdot10^{3}+3\cdot10^{2}+5\cdot10+2$
  \task $4\cdot10^{4}+7\cdot10^{2}+1\cdot10+8$
  \task $8\cdot10^{5}+5\cdot10^{3}+6\cdot10$
  \task $3\cdot10^{6}+4\cdot10^{2}+2$
\end{tasks}

\ex{2} \begin{tasks}(4)
  \task $72\,059$
  \task $300\,601$
  \task $4\,083\,000$
  \task $9009$
\end{tasks}

\ex{3} $5\cdot10^{6}+3\cdot10^{4}+7=\res{5\,030\,007}$.

\ex{3} $10^{9}\cdot10^{3}=10^{12}$: \ un 1 seguido de 12 ceros, $\res{1\,000\,000\,000\,000}$ operaciones.


% =====================================================
\seccio{7}{Propiedades de las potencias}

\ex{1} \begin{tasks}(6)
  \task $5^{6}$
  \task $7^{5}$
  \task $3^{8}$
  \task $2^{23}$
  \task $2^{10}$
  \task $3^{7}$
\end{tasks}

\ex{1} \begin{tasks}(6)
  \task $5^{6}$
  \task $3^{5}$
  \task $8^{3}$
  \task $2^{2}$
  \task $7^{3}$
  \task $4^{0}=1$
\end{tasks}

\ex{1} \begin{tasks}(6)
  \task $7^{6}$
  \task $5^{12}$
  \task $2^{10}$
  \task $9^{0}=1$
  \task $4^{12}$
  \task $6^{5}$
\end{tasks}

\ex{1} \begin{tasks}(6)
  \task $15^{4}$
  \task $14^{3}$
  \task $4^{5}$
  \task $3^{2}$
  \task $70^{3}$
  \task $6^{5}$
\end{tasks}

\ex{2} \begin{tasks}(3)
  \task $2^{8}:2^{4}=2^{4}$
  \task $5^{6}\cdot5^{3}=5^{9}$
  \task $6^{11}:6^{6}=6^{5}$
  \task $3^{18}:3^{10}=3^{8}$
  \task $3^{5}\cdot3^{2}=3^{7}$
  \task $7^{14}:7^{6}=7^{8}$
  \task $9^{4+3+14}=9^{21}$
  \task $4^{6}:4^{5}=4$
  \task $(3^{10})^{5}=3^{50}$
\end{tasks}

\ex{2} \begin{tasks}(3)
  \task $10^{5}=100\,000$
  \task $100^{3}=1\,000\,000$
  \task $3^{4}=81$
  \task $10^{3}=1000$
  \task $7^{2}=49$
  \task $1$
\end{tasks}

\ex{3} \begin{tasks}(5)
  \task $x=6$
  \task $x=2$
  \task $x=4$
  \task $x=2$
  \task $x=42$
  \task $x=8$
  \task $x=2$
  \task $x=0$
  \task $x=5$
  \task $x=1$
\end{tasks}

\ex{3} \begin{tasks}(1)
  \task $\left[2^{12}\cdot2^{5}\right]:2^{16}=2^{17}:2^{16}=\res{2}$
  \task $3^{5}\cdot3^{2}=\res{3^{7}}$
  \task No: $2^{3}+2^{3}=8+8=16=2^{4}$, pero $2^{6}=64$. Con las sumas no hay ninguna propiedad.
\end{tasks}


% =====================================================
\seccio{8}{Raíces cuadradas}

\ex{1} \begin{tasks}(4)
  \task $9$
  \task $8$
  \task $3$
  \task $10$
  \task $30$
  \task $13$
  \task $50$
  \task $11$
\end{tasks}

\ex{1} \begin{tasks}(4)
  \task $4$, resto $2$
  \task $6$, resto $4$
  \task $10$, resto $17$
  \task $3$, resto $6$
  \task $8$, resto $11$
  \task $5$, resto $6$
  \task $14$, resto $4$
  \task $9$, resto $16$
\end{tasks}

\ex{2} \begin{tasks}(4)
  \task $7$ y $8$
  \task $11$ y $12$
  \task $17$ y $18$
  \task $31$ y $32$
\end{tasks}\par Por ejemplo: $31^{2}=961<999<1024=32^{2}$.

\ex{2} \par\smallskip{\renewcommand{\arraystretch}{1.4}\centering\begin{tabular}{|C{3.2cm}|C{3.2cm}|C{3.2cm}|}\hline\rowcolor{soft}\textbf{Radicando} & \textbf{Raíz entera} & \textbf{Resto}\\\hline $69$ & $8$ & $5$\\\hline $144$ & $12$ & $0$\\\hline $245$ & $15$ & $20$\\\hline $60$ & $7$ & $11$\\\hline\end{tabular}\par}

\ex{3} Del $100$ al $120$ (el $121=11^{2}$ ya tiene raíz 11): resto máximo $120-100=\res{20}$ y hay $\res{21}$ números.

\ex{3} $15^{2}=225$, \ $16^{2}=256$, \ $17^{2}=289$, \ $18^{2}=324$, \ $19^{2}=361$.


% =====================================================
\seccio{9}{Operaciones y problemas con potencias y raíces}

\ex{1} \begin{tasks}(2)
  \task $9+8-8=\res{9}$
  \task $5+18-16=\res{7}$
  \task $6+16-12=\res{10}$
  \task $15+2=\res{17}$
  \task $8+3-9=\res{2}$
  \task $18+56=\res{74}$
\end{tasks}

\ex{2} \begin{tasks}(2)
  \task $5\cdot4+9=\res{29}$
  \task $27+6-24=\res{9}$
  \task $27-4+9\cdot16=\res{167}$
  \task $16+8-4=\res{20}$
  \task $3\cdot77=\res{231}$
  \task $4+20-1=\res{23}$
\end{tasks}

\ex{3} \begin{tasks}(1)
  \task $9\cdot100-8\cdot100=\res{100}$
  \task $2\cdot36+4\cdot400=72+1600=\res{1672}$
  \task $(1+14-9-5)\cdot(1+18-17)=1\cdot2=\res{2}$
  \task $\sqrt{100}-\sqrt{25}=10-5=\res{5}$
\end{tasks}

\ex{1} $36^{2}=\res{1296\text{ m}^{2}}$; \ $1296\cdot50=\res{64\,800\text{ €}}$.

\ex{1} $\sqrt{81}=\res{9\text{ sellos}}$ por lado.

\ex{2} Lado: $\sqrt{225}=15$ m. \ Perímetro: $4\cdot15=60$ m. \ Coste: $60\cdot6=\res{360\text{ €}}$.

\ex{2} $\sqrt{324}=\res{18\text{ filas}}$ (de 18 asientos).

\ex{2} $12^{2}=144\le150<169=13^{2}$: \ $\res{12}$ por lado, sobran $\res{6}$ y faltarían $169-150=\res{19}$.

\ex{3} Total: $5^{3}=\res{125}$. \ Sin pintar quedan los de dentro, un cubo de $3$ de arista: $3^{3}=\res{27}$.


\end{document}
